% Page budget: 1.55 pages | Owns: augmentation topology, additional states, encoder, and closed-loop objective.
% WRITING GUIDE
% Purpose: Define the final augmentation architecture and explain how it is trained inside the known feedback loop.
% Start here: Current implementation decisions D-129, D-140, D-141, D-180, D-220, D-233, and D-237 in docs/decisions.md.
% Supporting material: Quinten's 18 September outline feedback, the Meeting-audit synthesis, Thesis-writeup/Documentation/TRAINING-DESIGN.md, and docs/writeup figures.
% Mandatory papers: Read Hoekstra's 2025 LFR augmentation paper, the 2026 first-principles augmentation paper, the encoder-initialization paper, and Kessels' Chapter 5 before writing this section.
% Research-plan anchor: Preserve Aspect 2's dynamic parallel augmentation objective; document the departure from open-loop training and earlier state routing.
% Current decisions: Separate physical and additional states, keep the controller outside the model, score the complete training window without burn-in, and use one network for the physical correction and added-state update.
% Choices to justify: Final reported routing, added-state dimension, ANN width and depth, zero initialization, encoder history, closed-loop objective, rollout length, full-window scoring, and validation horizon. Use the configuration of the reported thesis run, not a superseded routing experiment.
% Horizon check: Treat encoder history, scored training window, joint-estimation horizon, and free-run validation horizon separately; relate candidates to relevant stable modes and test sensitivity to slower dynamics and free-integrator accumulation.
% Implementation audit: Read scripts/gantry/gantry_interconnect_dynamic.py, gantry_dynamic/{model,config,controller,data,training}.py, and model_augmentation/fit_systems/{interconnect,closed_loop,pre_encoder}.py.
% Reference check: Verify sources for parallel LFR augmentation, encoder initialization, closed-loop state extension, and any claimed architectural precedent against the original paper or thesis.
% Cautions: Earlier routing plans are superseded; distinguish the truth-only hidden absorber from learned states; closed-loop fit alone does not prove autonomous plant recovery.
% Figures: Absorber schematic, author's updated architecture, and thesis-owned common-reference loops. No residual panel or horizon figure.
% Section is finished when: Every signal has a defined frame and timing, the RK4 boundary is explicit, and the described architecture matches the final code.
%
% SOURCE MAP
% Hoekstra et al. 2025 EJC: dynamic-parallel topology, additional states, encoder-based truncated simulation.
% Hoekstra et al. 2026 Automatica submission: extended LFR formulation, baseline-preserving model/encoder initialization, joint identification context.
% Hoekstra et al. 2026 encoder paper: reconstructability-based W^b initialization; it does not identify W^a or test dynamic augmentation.
% Kessels 2025 thesis: direct closed-loop rollout, controller equations, and reconstruction of the machine controller state for each window.
% Drenth 2025 thesis: LPV-specific augmentation precedent; not the unrestricted MLP structure used here.
% This project: gantry specialization, residual-loop implementation, exact routing, RK4 boundary, and verification.
%
% MUST ESTABLISH (agreed with Dirk, 2026-10-06; step 2a of the thesis-section skill)
% Contribution delivered: the gantry realisation of S-DP augmentation (CT self-scheduled LPV-LFR baseline
%   inside a DT parallel model, joint estimation in the ten combinations) and its identification under
%   feedback. Not new: the S-DP topology (Hoekstra), augmentation of self-scheduled LPV-LFR with added
%   states (Drenth Eq. 5.2), closed-loop training with a known controller (Kessels, inspiration only).
% III-A Augmented model
%  1. Dynamic, not static: static augmentation adds no states (EJC Sec. 2, "flexible modes"); refitting
%     theta_base adds none either. Argue from the general case, no absorber.
%  1b. The augmented model is Hoekstra's LFR augmentation structure (EJC Eq. 4, 2026 Eq. 6) with a FIXED
%     interconnection (EJC Sec. 3.2) realising S-DP; its baseline block is Section II's self-scheduled
%     LFR with Delta(Y) = Y I6: an LFR inside an LFR, which is what the section title means.
%     Well-posedness: acyclic graph (EJC Thm. 1, 2026 Cond. 6) plus Section II's M(Y) condition.
%  2. Parallel, not series. LEAD REASON (Dirk, 2026-10-06): the orthogonal-by-construction method that
%     Section IV builds on is defined for the additive (parallel) structure (Gyorok 2026 abstract, Sec. 1;
%     D-003), so parallel lets us use that existing work. Supporting: Hoekstra gives no universal rule,
%     the structure follows prior knowledge (2026 Sec. 1, 6.3); the baseline stays the explicit nominal
%     term (research plan Aspect 2); parallel needs only a feedforward network for baseline-preserving
%     init, series needs ResNets (EJC Sec. 4.3, 2026 Sec. 5.4, 6.3).
%  3. State augmentation only, h_aug = 0: output P^T q is exact kinematics (cf. Kessels Remark 5.3). TODO reason.
%  4. Correction on all physical rows and added states: S-DP definition (EJC), X and Y reachable (D-103);
%     marginal-mode risk on free axes stated. TODO: why correct position rows rather than keep kinematics.
%     Context: Hoekstra 2026 Sec. 7 (F1Tenth) detaches the free integrators and learns only the
%     input-to-velocity map because the full state is measured; here only positions are measured.
%     One network for f_aug and g_aug is an implementation statement (D-180), no scientific claim.
%  5. DT correction after the CT RK4 step (input held, M(Y) per stage, affine normalisation). TODO reason
%     versus force injection in the CT derivative.
%  6. Unrestricted nonlinear augmentation (depends on Y through the state) versus LPV-structured
%     augmentation (Drenth): departure from the research plan's "in LPV-LFR form". TODO reason. Acyclic
%     routing adds no instantaneous loops; baseline well-posedness is Section II's.
%  7. Zero output layer reproduces the baseline from the same physical initial state (Hoekstra 2026
%     Eq. 29, D-237); the split is not unique (Sec. 5.2), handed to Section IV. Added states are memory
%     and model order, defined up to a state transformation, no physical coordinates.
% III-B Encoder
%  8. Encoder needed: only positions measured; truncated windows for cost and stability (Beintema).
%  9. W^b from the reconstructability map of a ZOH linearisation at an equilibrium (encoder paper
%     Eqs. 15-17, 30-32), not fitted (Hoekstra 2026 Eq. 30); W^a random. Claim only better starting
%     values and faster convergence, final accuracy comparable (encoder paper Sec. 4.4). The paper
%     assumes a stable baseline; ours is marginally stable (free X, Y), the map exists because the
%     positions make the linearisation observable. The paper treats co-estimating theta_base as out of
%     scope (footnote 1); we co-estimate. Precedent for linearising at zero: encoder paper Sec. 4.2
%     ("around the 0 point"), no reason given. TODOs: why Y = 0; disclose nominal parameters in W^b
%     versus detuned start; why model-based rather than data-based init (candidate: cost, Sec. 4.4);
%     Kaiming versus Xavier (todo only).
% 10. Encoder estimates the full state incl. measured positions (SUBNET). TODO: mention Kessels Remark 5.3?
% 11. Keep encoder history, scored window and validation horizon apart; values in Setup. The
%     joint-estimation sensitivity horizon belongs to Section IV.
% III-C Closed-loop identification
% 12. Known-controller, reference-driven output-error identification, inspired by Kessels (Eqs. 5.12,
%     5.13; Remark 5.6 open-loop attempt failed). Related to the indirect approach (Forssell and Ljung
%     Sec. 5.2) as a linear-theory motivation, not a guarantee; direct identification is not called
%     inherently biased. Free X and Y axes: the SUBNET consistency argument that Hoekstra inherits
%     (2026 Sec. 5.5) assumes an incrementally exponentially output stable data-generating system
%     (Beintema 2023 Condition 1); the open-loop gantry is not (free integrators), the loop with the
%     stabilising controller can be. Precondition, not a proof, for our nonlinear model. Same principle
%     as Kessels (closed-loop simulation of an unstable open-loop system), scoped to marginal stability.
% 13. Our formulation (ours, "we"): u_hat = u + C x^c + D(y - y_hat). Derivation by subtracting the
%     machine's controller equations, equal to the shared-r, shared-u_ff loop of the figure under a
%     compatible controller, signals and initial state (assumption; whether the data meet it is a Setup
%     or Results check: 4 kHz controller on recorded r - y versus recorded feedback force).
%     Main consequence: x^c_tau = 0 per window, no controller initialisation; consistent with the
%     encoder (model state and model controller both start on the machine's condition at tau). Contrast:
%     Kessels' direct form needs the machine controller state (Remark 5.4). Error before tau is not
%     carried in; the free-run validation does carry it.
% 14. Step order from no plant feedthrough: no algebraic loop, no added delay.
% 15. Every window sample scored (as Eq. 22); theta_base, theta_aug, eta joint, controller fixed.
%     Joint estimation framing (Dirk, 2026-10-06): the goal is to estimate theta_base jointly, in the ten
%     combinations. Obstacle from Jan's work: the split is not unique (2026 Sec. 5.2) and approximately
%     initialised parameters stay near their start (EJC Sec. 4.3, 2026 Sec. 6.3). Handover: Section IV
%     extends Jan's method with orthogonality so that the estimate of theta_base is UNIQUE (this
%     strength, not "recovers the true parameters"). State what interpretable means: the estimated
%     parameters keep their physical meaning; cite Gyorok 2026 Sec. 3 ("With non-unique theta_b, the
%     baseline model could lose its physical meaning"). Operational form only; the general definition
%     is the Introduction's (todo there). "Negation" and Kessels' definition (Sec. 5.3.1.3) belong to
%     Section IV; his Sec. 6.2 recommendation to the Introduction gap. Do NOT mention Jan's parameter regularisation anywhere.
%     Do NOT say explicitly that Section III's formulation is the unconstrained comparison arm; let
%     Section IV extend it naturally.
% Setup pattern for hyperparameters (2026 Sec. 7): honest basis per value (physical insight and
%     trial-and-error, line-search, inherited from earlier results).
% 16. Closed-loop fit can mask plant error; plant-domain checks and controller transfer in Results.
% Notation: u = P u_act as in Sections II and IV; no "(eom)" shorthand; controller matrices distinct from A_c(Y).
%
% MUST ESTABLISH, PROPOSED (proposal-test-III, 2026-10-09; no approval available; inferred from the README
%   ownership row III, its Math standard with the PS2 example, and source-map rows III-A to III-D):
% Structure: the skill's Order rules give the closed-loop simulation its own subsection, so items 12 to 14
%   form III-C (sec:aug_loop, the label Sections II and IV already cite) and items 15, 16 with the training
%   coordinates, optimiser, selection rule and the name M_U form III-D (sec:aug_closed_loop). PS2 is III-E
%   (sec:aug_ps2), after the criterion it opens (Hoekstra 2026 Sec. 5: criterion before initialisation).
%   D-243 calls PS2 "III-D"; Section IV cites labels only.
% III-E Added-state initialisation (PS2; this work, adapted from Hefny 2015 Sec. 2 and Downey 2017 Sec. 4.2):
% 17. Problem: from the zero start the added states are zero after the first step and V gives W^a and the
%     added-state output rows zero gradient; plain training can lower V through f_aug alone (static
%     corrections, real lags; PS2-METHOD Sec. 2.2); whether it does is THESIS-RESULTS step 2, stated no stronger.
% 18. S1 display: the encoder's added-state outputs and a temporary decoder fitted to the current model's
%     normalised closed-loop residual over M samples; trains theta_p and theta^a_enc only, own optimiser, one
%     update after every Adam update of V, calibration records withheld; fallback when unpredictable.
% 19. S2 display: once, one-step regression of g_aug on frozen encoder pairs (k, k+1); only the added-state
%     rows enter the loss, the shared hidden layers move.
% 20. S3: decoder discarded, eq:aug_loss continues unchanged; deployed model, criterion and selection
%     unchanged; added states stay latent; M_PS2 named.
% Notation set in this session: encoder window chi_tau (xi is the training coordinate Section IV uses);
%   encoder parameters theta_enc and its added-state part theta^a_enc (Section IV); machine controller state
%   x^fb (s is the m_Delta scale of the appendix); normalised signals carry a superscript n (Section IV);
%   the model input is u_act, as in the code and Section II-D, with P inside f_base through eq:eom.
\section{Dynamic LPV-LFR Augmentation}\label{sec:augmentation}
\ifdraft
\begin{itemize}
\item The section extends the baseline with learned dynamics and identifies the result from closed-loop records; it names its five subsections (style profile)
\item Contribution share: the realisation on the self-scheduled gantry baseline with joint estimation of $\theta_{\mathrm{base}}$, its identification in closed loop with the known controller, and the PS2 initialisation of the added states (MUST ESTABLISH)
\end{itemize}
\fi

Given the baseline of Section~\ref{sec:lfr}, our goal is a model that
also captures the dynamics the baseline omits, identified from records
measured under feedback.
To this end, we define the augmented model (Section~\ref{sec:aug_structure}),
the encoder that estimates its initial state (Section~\ref{sec:aug_encoder})
and its simulation inside the known feedback loop
(Section~\ref{sec:aug_loop}).
We then state the identification problem
(Section~\ref{sec:aug_closed_loop}) and an initialisation of the added states
that opens its optimisation (Section~\ref{sec:aug_ps2}).
This delivers the second contribution of Section~I, dynamic parallel LFR
augmentation of the self-scheduled gantry baseline.
Its parameters $\theta_{\mathrm{base}}$ are estimated jointly in the
combinations of Section~\ref{sec:params}, and the model is identified in
closed loop with the known controller.
\todo{Missing: the Introduction's contribution list does not name PS2, which this section adds as this work's method. Candidate: extend the second contribution with ``and a data-driven initialisation of the added states'' (basis: D-236, D-243, \texttt{THESIS-RESULTS.md} step~3).}

\subsection{Augmented model}\label{sec:aug_structure}
\ifdraft
\begin{itemize}
\item If the omitted effects are dynamic, the six baseline states cannot represent them; refitting $\theta_{\mathrm{base}}$ or a static augmentation adds no states (EJC Sec.~2)
\item The model is Hoekstra's LFR augmentation structure with a fixed interconnection realising S-DP; its baseline block is the LFR of Section~II (EJC Sec.~3, 2026 Table~1)
\item Parallel because the orthogonal-by-construction method of Section~IV is defined for the additive structure (Gy\"or\"ok 2026); the baseline stays an explicit term
\item State augmentation only: the output map is the exact kinematic relation (Kessels Remarks~5.3, 5.1)
\item $f_{\mathrm{base}}$ is one RK4 step with the input held and $M(Y)$ per stage, so it stays self-scheduled; affine normalisation from the training records, unlike Hoekstra Eq.~(28)
\item The learned correction is a discrete-time correction after the step (?)
\item One feedforward network with tanh hidden layers supplies $f_{\mathrm{aug}}$ and $g_{\mathrm{aug}}$; plain feedforward suffices for a parallel structure (EJC Sec.~4.3)
\item The correction acts on all physical rows, $X$ and $Y$ included; on these axes it meets no stiffness (?: why the position rows)
\item The network depends on $Y$ through the state but has no LPV structure, unlike Drenth Eq.~(5.2) (?)
\item The interconnection is acyclic, so it adds no algebraic loop; the baseline loop is well posed by Section~II
\item A zero output layer reproduces the baseline from the same physical initial state (2026 Eq.~(29)); what this does to the added states is Section~III-E's problem
\item The baseline/network split is not unique (2026 Sec.~5.2); interpretable means the estimated $\theta_{\mathrm{base}}$ keeps its physical meaning (Gy\"or\"ok Sec.~3); Section~IV makes the estimate unique
\item Normalisation: affine, statistics of the training records, unlike the scaling of 2026 Eq.~(28); later displays in normalised symbols (Math standard rule~4)
\end{itemize}
\fi

If the effects that the baseline omits are dynamic, its six states cannot
represent them.
Refitting $\theta_{\mathrm{base}}$ adds no states.
Neither does a static augmentation, which corrects the state update of the
baseline without extending its state~\cite[Sec.~2]{hoekstra2025lfr}.
The augmentation therefore needs states of its own.

We use the LFR-based augmentation structure of Hoekstra et al., in which an
interconnection matrix connects a baseline block and a learning
block~\cite[Sec.~3.1]{hoekstra2025lfr}, here a fixed
one~\cite[Sec.~3.2]{hoekstra2025lfr}.
Its baseline block is the self-scheduled LFR of Section~\ref{sec:lfr}, so the
augmented model is an LFR whose baseline block is itself an LFR in $Y$.
Of the interconnections, we choose the dynamic-parallel
one~\cite[Table~1]{hoekstra2026lfr}, because the orthogonal-by-construction
parametrisation that Section~\ref{sec:obc} extends is defined for this
additive structure~\cite{gyorok2026obc}.
It also keeps the baseline an explicit term of the state update.
The output map is not augmented, because it is the exact kinematic relation
\eqref{eq:Ptransform}.
Kessels likewise leaves an accurate output map
unaugmented~\cite[Remark~5.3]{kessels2025ai}.
An output correction that depends on the input would also make the output
depend on the feedback force computed from it~\cite[Remark~5.1]{kessels2025ai}.

The baseline transition $f_{\mathrm{base}}$ is one fourth-order Runge-Kutta
(RK4) step of \eqref{eq:eom} over the sample period $T_s$, with the input held
over the step.
Each stage evaluates $M(Y)$ at its own state, so $f_{\mathrm{base}}$ remains
self-scheduled.
The learned correction is added to the result of the step, so it is a
discrete-time correction at the sample rate, while $f_{\mathrm{base}}$
integrates the continuous-time physics.
In physical units, the augmented model is
\begin{equation}
\begin{aligned}
 \tilde x_{k+1}
 &=f_{\mathrm{base}}(\theta_{\mathrm{base}},\tilde x_k,u_{\mathrm{act},k})
   +f_{\mathrm{aug}}(\theta_{\mathrm{aug}},\hat x_k,u_{\mathrm{act},k}),\\
 \bar x_{k+1}
 &=g_{\mathrm{aug}}(\theta_{\mathrm{aug}},\hat x_k,u_{\mathrm{act},k}),\\
 \hat y_k&=P^\top q_k,
\end{aligned}
\label{eq:aug_transition}
\end{equation}
where $\tilde x_k=\col(q_k,\dot q_k)\in\R^6$ is the physical state of
\eqref{eq:lfr_G}, $\bar x_k\in\R^{n_{\bar x}}$ the added state,
$\hat x_k=\col(\tilde x_k,\bar x_k)$, $u_{\mathrm{act},k}\in\R^3$ the actuator
forces, $\hat y_k\in\R^3$ the predicted actuator positions, whose recorded
counterpart $q_{\mathrm{act},k}$ we write $y_k$,
$\theta_{\mathrm{base}}\in\R^{10}$ the combinations \eqref{eq:phi},
$\theta_{\mathrm{aug}}\in\R^{n_{\mathrm{aug}}}$ the network parameters, and
$f_{\mathrm{base}}:\R^{10}\times\R^6\times\R^3\to\R^6$,
$f_{\mathrm{aug}}:\R^{n_{\mathrm{aug}}}\times\R^{6+n_{\bar x}}\times\R^3\to\R^6$
and $g_{\mathrm{aug}}:\R^{n_{\mathrm{aug}}}\times\R^{6+n_{\bar x}}\times\R^3\to\R^{n_{\bar x}}$.
No learned function feeds the baseline block or itself within a step, so the
interconnection graph is acyclic~\cite[Thm.~1]{hoekstra2025lfr} and, with
the well-posedness of the baseline LFR (Section~\ref{sec:lfr}), the augmented
model is well posed for every $Y$.
\todo{Missing: reason for correcting the state after the RK4 step rather than adding a learned force inside the continuous-time derivative. No decision entry gives a reason other than gradient magnitude.}

The network and the identification act on centred signals of unit variance.
Hoekstra et al.\ only scale the signals and take the state scale from a
simulation of the baseline~\cite[Sec.~5.3]{hoekstra2026lfr}.
We also centre them and take every statistic from the training records, the
state statistics from the positions $q=P^{-\top}y$ and their first
differences:
\begin{equation}
 \tilde x^{\mathrm n}=T_x(\tilde x-\mu_x),\quad
 u^{\mathrm n}=T_u(u_{\mathrm{act}}-\mu_u),\quad
 y^{\mathrm n}=T_y(y-\mu_y),
 \label{eq:aug_norm}
\end{equation}
where $\mu_x\in\R^6$, $\mu_u,\mu_y\in\R^3$ are the channel means,
$T_\bullet=\operatorname{diag}(\sigma_\bullet)^{-1}$ with $\sigma_\bullet$ the
channel standard deviations, and $\bar x$ is not scaled, so
$\hat x^{\mathrm n}=\col(\tilde x^{\mathrm n},\bar x)$.
We write $f^{\mathrm n}_{\mathrm{base}}$, $f^{\mathrm n}_{\mathrm{aug}}$ and
$g^{\mathrm n}_{\mathrm{aug}}$ for the maps of \eqref{eq:aug_transition} in
these coordinates.
Since the position part of $\mu_x$ is $P^{-\top}\mu_y$, the output map stays
linear,
$\hat y^{\mathrm n}_k=C^{\mathrm n}\tilde x^{\mathrm n}_k$ with
$C^{\mathrm n}=T_yP^\top[I_3\;\,0]\,T_x^{-1}$.

One multilayer perceptron with tanh hidden layers and a linear output layer
supplies both learned functions.
A feedforward network suffices because the baseline is augmented additively,
whereas a series structure needs a residual network to reproduce the baseline
at initialisation~\cite[Sec.~4.3]{hoekstra2025lfr}.
Its first six outputs correct every physical row, as in the dynamic-parallel
structure~\cite[Sec.~2]{hoekstra2025lfr}, so the learned component also acts
on the free axes $X$ and $Y$, which have no stiffness in
\eqref{eq:damping_stiffness_matrices}.
The network reads $Y$ through $\hat x^{\mathrm n}$ but, unlike the augmented
LPV-LFR of Drenth~\cite[Eq.~(5.2)]{drenth2025thesis}, has no LPV structure:
\begin{equation}
 \begin{bmatrix}
  f^{\mathrm n}_{\mathrm{aug}}(\theta_{\mathrm{aug}},\hat x^{\mathrm n}_k,u^{\mathrm n}_k)\\
  g^{\mathrm n}_{\mathrm{aug}}(\theta_{\mathrm{aug}},\hat x^{\mathrm n}_k,u^{\mathrm n}_k)
 \end{bmatrix}
 =\phi_{\mathrm{aug}}(\theta_{\mathrm{aug}},\hat x^{\mathrm n}_k,u^{\mathrm n}_k),
 \label{eq:aug_network}
\end{equation}
where
$\phi_{\mathrm{aug}}:\R^{n_{\mathrm{aug}}}\times\R^{6+n_{\bar x}}\times\R^3\to\R^{6+n_{\bar x}}$
(Fig.~\ref{fig:aug_structure}).
A persistent correction on the $X$ or $Y$ row therefore accumulates in
position.
\todo{Missing: why the position rows are corrected rather than kept as the integrals of the velocities. Context: Hoekstra et al.\ detach the free integrators of their vehicle and learn only the input-to-velocity map, which needs measured velocities~\cite[Secs.~7.2, 7.4]{hoekstra2026lfr}; here only positions are measured. D-103 requires only that $X$ and $Y$ are reachable.}
\todo{Missing: reason for an unrestricted network rather than an LPV-structured augmentation. The research plan (Aspect~2) proposed the interconnection in LPV-LFR form; this is a departure, and D-017 is not superseded by a later decision.}

\begin{figure}[t]
 \centering
 \resizebox{\columnwidth}{!}{\begin{tikzpicture}[
    >={Latex[length=1.8mm,width=1.3mm]},
    line width=0.5pt,
    blk/.style={draw, line width=0.6pt, rounded corners=1pt, fill=white,
                align=center, inner sep=2pt, minimum height=9mm},
    sub/.style={font=\scriptsize},
    dot/.style={circle, fill, inner sep=0pt, minimum size=2.4pt},
    hop/.style={preaction={draw=white, line width=2.4pt}},
    sum/.style={draw, circle, line width=0.6pt, inner sep=0pt, minimum size=3.4mm},
  ]
  % grid
  \def\xU{0}     % u_k bus
  \def\xS{0.8}   % state bus x-hat_k
  \def\xB{3.25}  % block centres
  \def\xP{5.35}  % sums
  \def\xR{6.55}  % merge bar
  \def\xF{7.45}  % feedback
  \def\yA{2.6}   % phi_aug row
  \def\yF{1.3}   % f_base row
  \def\yO{0.0}   % h_base row
  \def\yR{-1.25} % z^-1
  \def\yE{-2.65} % encoder, below z^-1

  % blocks
  \node[blk, minimum width=27mm] (ann) at (\xB,\yA)
       {$\phi_{\mathrm{aug}}$\\[-1pt] {\scriptsize neural network}};
  \node[blk, minimum width=27mm] (fb) at (\xB,\yF)
       {$f_{\mathrm{base}}$\\[-1pt] {\scriptsize physics model, one step}};
  \node[blk, minimum width=27mm] (hb) at (\xB,\yO)
       {$h_{\mathrm{base}}$\\[-1pt] {\scriptsize output map}};
  \node[blk, minimum width=8mm, minimum height=7mm] (z) at (\xS,\yR) {$z^{-1}$};
  \node[blk, minimum width=11mm, minimum height=11mm] (enc) at (\xS,\yE)
       {$\psi$\\[-1pt] {\scriptsize encoder}};

  \node[sum] (s1) at (\xP,\yF) {};
  \node[sum] (s2) at (\xP,\yO) {};
  \foreach \s in {s1,s2}
    \draw ($(\s)+(-0.95mm,0)$) -- ($(\s)+(0.95mm,0)$)
          ($(\s)+(0,-0.95mm)$) -- ($(\s)+(0,0.95mm)$);

  % state bus x-hat_k
  \draw (z.north) -- (\xS,\yA+0.2);
  \node[anchor=east] at (\xS,{(\yR+\yO)/2+0.05}) {$\hat x_k$};
  \draw[->] (\xS,\yA+0.2) -- (ann.west |- 0,\yA+0.2);
  \draw[->] (\xS,\yF+0.2) -- (fb.west |- 0,\yF+0.2)
       node[pos=0.45, above, inner sep=1.5pt] {$\tilde x_k$};
  \draw[->] (\xS,\yO) -- (hb.west)
       node[pos=0.45, above, inner sep=1.5pt] {$\tilde x_k$};
  \node[dot] at (\xS,\yF+0.2) {};
  \node[dot] at (\xS,\yO) {};

  % input bus u_k (crosses the state bus with a gap)
  \node[anchor=south, inner sep=1.5pt] at (\xU,\yA+0.7) {$u_k$};
  \draw (\xU,\yA+0.7) -- (\xU,\yF-0.2);
  \draw[->, hop] (\xU,\yA-0.2) -- (ann.west |- 0,\yA-0.2);
  \draw[->, hop] (\xU,\yF-0.2) -- (fb.west |- 0,\yF-0.2);
  \node[dot] at (\xU,\yA-0.2) {};

  % network outputs
  \draw[->] (ann.east |- 0,\yA+0.2) -- (\xR,\yA+0.2)
       node[pos=0.22, above, inner sep=1.5pt] {$g_{\mathrm{aug}}$}
       node[pos=0.8, above, inner sep=1.5pt] {$\bar x_{k+1}$};
  \draw[->] (ann.east |- 0,\yA-0.2) -| (s1.north)
       node[pos=0.75, right, inner sep=1.5pt] {$f_{\mathrm{aug}}$};

  % physics step -> sum -> merge bar
  \draw[->] (fb.east) -- (s1.west);
  \draw[->] (s1.east) -- (\xR,\yF)
       node[pos=0.5, above, inner sep=1.5pt] {$\tilde x_{k+1}$};
  \draw[line width=2pt] (\xR,\yF-0.2) -- (\xR,\yA+0.4);

  % feedback x-hat_{k+1} -> z^-1
  \draw[->] (\xR,{(\yA+\yF)/2}) -- (\xF,{(\yA+\yF)/2}) |- (z.east);
  \node[anchor=south, inner sep=1.5pt] at ({(\xR+\xF)/2+0.03},{(\yA+\yF)/2}) {$\hat x_{k+1}$};

  % output
  \draw[->] (hb.east) -- (s2.west);
  \draw[->, dashed] (\xP,\yO+0.7) -- (s2.north);
  \node[anchor=west, inner sep=1.5pt] at (\xP+0.05,\yO+0.6) {$h_{\mathrm{aug}}=0$};
  \draw[->] (s2.east) -- ++(0.55,0) node[anchor=west, inner sep=1.5pt] {$\hat y_k$};

  % encoder
  \draw[->, dashed] (enc.north) -- (z.south)
       node[pos=0.5, right, inner sep=1.5pt] {$\hat x_{k|k}$};
  \draw[<-] (enc.west |- 0,\yE+0.3) -- ++(-0.4,0)
       node[anchor=east, inner sep=1.5pt] {$y_{k-n_a}^{k-1}$};
  \draw[<-] (enc.west |- 0,\yE-0.3) -- ++(-0.4,0)
       node[anchor=east, inner sep=1.5pt] {$u_{k-n_b}^{k-1}$};
\end{tikzpicture}}
 \caption{Augmented model. The physics step $f_{\mathrm{base}}$ and the
 network $\phi_{\mathrm{aug}}$ act in parallel on $\hat x_k$.
 $\phi_{\mathrm{aug}}$ supplies $f_{\mathrm{aug}}$ and $g_{\mathrm{aug}}$,
 and the output map is not learned. The encoder
 $\psi$ sets $\hat x$ at each window start. During training, the input is the
 closed-loop input of Section~\ref{sec:aug_loop}.
 \todo{No setpoint generator currently for the inputs. to cramped together. We could make this a two column figure.}
 }
 \label{fig:aug_structure}
\end{figure}

Training starts from the baseline, as Hoekstra et al.\ require of the
initialisation~\cite[Eq.~(29)]{hoekstra2026lfr}.
The output layer of $\phi_{\mathrm{aug}}$ is set to zero,
\begin{equation}
 W_L=0,\quad b_L=0
 \quad\Longrightarrow\quad
 f^{\mathrm n}_{\mathrm{aug}}=0,\quad g^{\mathrm n}_{\mathrm{aug}}=0,
 \label{eq:aug_zero_init}
\end{equation}
where $W_L\in\R^{(6+n_{\bar x})\times n_h}$ and $b_L\in\R^{6+n_{\bar x}}$ are
its weight and bias and $n_h$ is the hidden width.
From the same physical initial state, the augmented model then reproduces the
baseline.
\todo{Missing: attribution of the all-zero output layer. It follows the authors' public dynamic-parallel code (D-237) and is the thesis choice (D-240); \cite[Sec.~5.4.3]{hoekstra2026lfr} also allows a random linear part. Ask Jan which produced the reported S-DP results.}

From this start, $f^{\mathrm n}_{\mathrm{aug}}$ can also learn changes that a
different $\theta_{\mathrm{base}}$ would produce, so the split between
baseline and network is not unique~\cite[Sec.~5.2]{hoekstra2026lfr}.
In the experiments of Hoekstra et al., approximately initialised physical
parameters stay close to their initial values, while the network learns
dynamics the baseline could represent~\cite[Sec.~6.3]{hoekstra2026lfr}.
With a non-unique $\theta_{\mathrm{base}}$, the baseline can lose its physical
meaning~\cite[Sec.~3]{gyorok2026obc}.
Here, the estimated combinations are interpretable if they keep that meaning,
which requires a unique estimate.
Section~\ref{sec:obc} extends the method so that this estimate is unique.

\subsection{Encoder and initial state}\label{sec:aug_encoder}
\ifdraft
\begin{itemize}
\item Each window starts from a state of which only the positions are measured; an encoder estimates it from the input-output history (Beintema)
\item Windows starting from an estimated state keep the cost independent of the record length and stabilise the optimisation (Beintema Secs.~2, 3); the window includes sample $\tau$, unlike 2026 Eq.~(22e) (encoder paper Eqs.~(11) to (17))
\item The encoder estimates all states, the measured positions included (?: Kessels Remark~5.3)
\item The encoder is a linear map plus a nonlinear correction (encoder paper Eq.~(8)); the linear part acts on scaled, uncentred histories (D-055, code)
\item $W^a$ starts from a Xavier draw (2026 Eq.~(31); D-239; adopted); the correction starts at zero
\item $W^b$ starts from the reconstructability map of the baseline linearised at rest, $Y=0$ (encoder paper Eqs.~(15) to (17), (30) to (32)); computed, not fitted as in 2026 Eq.~(30)
\item The paper assumes a stable baseline; ours is marginally stable, but the positions make the linearisation observable
\item The initialisation improves starting values and convergence speed, not final accuracy (encoder paper Sec.~4.4)
\item Open: why $Y=0$, nominal parameters in $W^b$, model-based versus data-based initialisation (?)
\end{itemize}
\fi

Only the positions of the state at the start $\tau$ of a training window are
measured.
Following Beintema et al., an encoder estimates this state from the input and
output history~\cite{beintema2023subnet}.
Simulating short windows from encoded states keeps the cost of a gradient step
independent of the record length and stabilises the
optimisation~\cite[Secs.~2, 3]{beintema2023subnet}.
Unlike the criterion of Hoekstra et al., whose encoder reads the samples up to
$\tau-1$~\cite[Eq.~(22e)]{hoekstra2026lfr}, the history includes sample
$\tau$, as the reconstructability map used below
requires~\cite[Eqs.~(11) to (17)]{hoekstra2026encoder}:
\begin{equation}
\begin{aligned}
 \hat x^{\mathrm n}_{\tau|\tau}
 &=\psi(\theta_{\mathrm{enc}},\chi_\tau),\\
 \chi_\tau&=\col\bigl(y^{\mathrm n}_{\tau-n},\dots,y^{\mathrm n}_\tau,\,
   u^{\mathrm n}_{\tau-n},\dots,u^{\mathrm n}_\tau\bigr),
\end{aligned}
 \label{eq:aug_encoder}
\end{equation}
where $n$ is the encoder history, $\chi_\tau\in\R^{6(n+1)}$ the recorded
window, $\theta_{\mathrm{enc}}\in\R^{n_{\mathrm{enc}}}$ the encoder parameters
and $\psi:\R^{n_{\mathrm{enc}}}\times\R^{6(n+1)}\to\R^{6+n_{\bar x}}$.
The encoder thus estimates all states, the measured positions included.
\todo{Missing: whether to mention the alternative of initialising the positions from the measurement~\cite[Remark~5.3]{kessels2025ai}.}
\todo{Align the figure's past-only encoder-history labels with
\eqref{eq:aug_encoder}, which also uses the current sample.}

The encoder is a linear map with a nonlinear
correction~\cite[Eq.~(8)]{hoekstra2026encoder}.
Its added-state rows $W^a$ have no baseline counterpart and start from a
Xavier draw, as in~\cite[Eq.~(31)]{hoekstra2026lfr}.
The correction starts at zero.
The linear part acts on the scaled but uncentred window, for which the map
\eqref{eq:aug_wb} below is built:
\begin{equation}
 \psi(\theta_{\mathrm{enc}},\chi_\tau)
 =\begin{bmatrix}W^b\\ W^a\end{bmatrix}(\chi_\tau+\chi_0)
  -\begin{bmatrix}T_x\mu_x\\ 0\end{bmatrix}
  +\tilde\psi(\chi_\tau),
 \label{eq:aug_enc_lin}
\end{equation}
where $W^b\in\R^{6\times6(n+1)}$, $W^a\in\R^{n_{\bar x}\times6(n+1)}$,
$\chi_0$ restores the means that \eqref{eq:aug_norm} removes, $\tilde\psi$ is
a multilayer perceptron whose output layer starts at zero, and
$\theta_{\mathrm{enc}}$ collects $W^b$, $W^a$ and the parameters of
$\tilde\psi$.
We write $\psi_a$ for the last $n_{\bar x}$ rows of $\psi$ and
$\theta^a_{\mathrm{enc}}$ for $W^a$ with the added-state rows of the output
layer of $\tilde\psi$.
\todo{VERIFY bib: \texttt{hoekstra2026encoder} = arXiv:2602.13108v1
(13 Feb 2026), Hoekstra, Gy\"or\"ok, T\'oth, Schoukens; matches PDF p.~1.
Check for a published (IFAC) version before submission.}

The physical rows start from the baseline, through the model-based
initialisation of Hoekstra et al.~\cite[Sec.~3.1]{hoekstra2026encoder}.
Its map needs a linear model, so as in~\cite[Eqs.~(30) to (32)]{hoekstra2026encoder}
we linearise the baseline at an equilibrium: at rest at $Y=0$, which is one
for zero input because $K$ in \eqref{eq:damping_stiffness_matrices} vanishes
on $X$ and $Y$.
Discretising by zero-order hold over $T_s$ with the nominal parameters gives
$(A_d,B_d,C_d)$ in the scaled coordinates.
The paper assumes a stable baseline~\cite[Sec.~2]{hoekstra2026encoder}.
Ours is marginally stable, but the measured positions make $(A_d,C_d)$
observable, so the observability matrix has full column rank.
The paper also leaves co-estimation of the baseline parameters out of
scope~\cite[footnote~1]{hoekstra2026encoder}, whereas here $W^b$ is trained
jointly with them.
Unlike~\cite[Eq.~(30)]{hoekstra2026lfr}, which fits the baseline encoder to
simulated baseline states, we compute the reconstructability
map~\cite[Eqs.~(15) to (17)]{hoekstra2026encoder} directly:
\begin{equation}
 W^b=\begin{bmatrix}
  A_d^{\,n}O_n^{+} & \;r_n-A_d^{\,n}O_n^{+}\Gamma_n
 \end{bmatrix},
 \label{eq:aug_wb}
\end{equation}
where $O_n\in\R^{3(n+1)\times6}$ is the observability matrix of $(A_d,C_d)$
over $n+1$ samples, $\Gamma_n$ the Toeplitz matrix of its Markov parameters,
$r_n$ the input-to-state map over the same samples, ${}^{+}$ the Moore-Penrose
pseudoinverse, and the two blocks act on the output and input parts of
$\chi_\tau$.
This start gives better starting values and faster convergence than a random
encoder, at comparable final accuracy~\cite[Sec.~4.4]{hoekstra2026encoder}.
\todo{Missing: why linearise at $Y=0$. The encoder paper also linearises ``around the 0 point''~\cite[Sec.~4.2]{hoekstra2026encoder} but gives no reason.}
\todo{Missing: disclosure. $W^b$ is built from the nominal parameters, which equal those of the simulated truth (Section~\ref{sec:setup}), while $\theta_{\mathrm{base}}$ may start detuned and is co-estimated, which the encoder paper leaves out of scope~\cite[footnote~1]{hoekstra2026encoder}. Candidate: state it as prior information the detuned runs receive, or build $W^b$ at the start $\theta_{\mathrm{base},0}$ (basis: \texttt{GD/model.py::build\_model} calls \texttt{gantry\_linearize\_and\_discretize} with the nominal constants).}
\todo{Missing: why the model-based rather than the data-based initialisation~\cite[Sec.~3.2]{hoekstra2026encoder}. Candidate: negligible computation~\cite[Sec.~4.4]{hoekstra2026encoder}.}

\subsection{Closed-loop simulation}\label{sec:aug_loop}
\ifdraft
\begin{itemize}
\item The records are measured under feedback and the model must predict the machine under feedback, so each window is simulated in the loop, not driven by the recorded input (2026 Eq.~(22))
\item Inspired by Kessels (Eqs.~(5.12), (5.13); Remark~5.6); for linear systems the indirect approach, where an output-error criterion remains consistent (Forssell and Ljung); a motivation, not a guarantee
\item $X$ and $Y$ are free, so the open-loop gantry violates the stability condition behind SUBNET consistency (Beintema Cond.~1); the loop with the controller can satisfy it
\item Both loops share $r$, $u^{\mathrm{ff}}$ and the controller; subtracting them gives our residual form, exact under a compatible controller, signals and initial state (check in Setup)
\item $x^{c}_\tau=0$ needs no controller initialisation and matches the encoder; Kessels' direct form needs the machine's controller state (Remark~5.4)
\item No plant feedthrough fixes the step order: no algebraic loop, no added delay; error before $\tau$ is not carried into the window
\end{itemize}
\fi

The records are measured under feedback, and the model must predict the
machine under feedback: its response to a reference, a feedforward and a
known controller.
We therefore simulate each training window in closed loop with that
controller, rather than driving the model with the recorded input as
in~\cite[Eq.~(22)]{hoekstra2026lfr}.
Kessels trains augmented models in this way, with the known controller around
the model~\cite[Eqs.~(5.12), (5.13)]{kessels2025ai}.
He reports that identifying the open-loop model from his closed-loop data
failed and attributes this to noise-induced
bias~\cite[Remark~5.6]{kessels2025ai}.
Our records are noisy as well (Section~\ref{sec:setup}).
For linear systems, fitting the closed loop from the reference with a known
controller is the indirect approach to closed-loop identification, in which
an output-error criterion remains consistent because the reference is
uncorrelated with the noise~\cite{forssell1999revisited}.
For our nonlinear model, this is a motivation, not a guarantee.
\todo{Missing: D-140 (Ruled out~(4)) barred Remark~5.6 as an argument while the records were noise-free; the thesis records are noisy (D-225). Candidate: the ban no longer holds (basis: D-225, \texttt{THESIS-RESULTS.md} Sec.~2, noisy training records).}
\todo{Missing: section number for~\cite{forssell1999revisited}. The indirect approach was verified in Sec.~5.2.1 of the LiU report LiTH-ISY-R-2021; the Automatica numbering is unchecked. Candidate: cite without a section number until the journal version is read.}

The free axes of the baseline give a second reason.
$X$ and $Y$ have no stiffness in \eqref{eq:damping_stiffness_matrices}, so the
open-loop gantry is marginally stable: a difference in initial velocity leaves
a permanent position offset.
The consistency argument of the identification
method~\cite[Sec.~5.5]{hoekstra2026lfr} rests on that of SUBNET, which assumes
an incrementally exponentially output stable data-generating
system~\cite[Cond.~1]{beintema2023subnet}.
The open-loop gantry does not satisfy this condition, whereas the loop with
the stabilising controller can.
Kessels uses the same principle to update models of systems that are unstable
in open loop~\cite[Ch.~5]{kessels2025ai}.
For our nonlinear model, this makes the condition attainable but does not
prove consistency.
\todo{Missing: evidence on the thesis data that open-loop rollouts accumulate position error on $X$ and $Y$. Candidate: the open-loop versus closed-loop drift on one record that \texttt{THESIS-RESULTS.md} Sec.~5 plans, reported in Section~\ref{sec:results}.}

Both loops are driven by the same reference $r$ and feedforward
$u^{\mathrm{ff}}$ through the same controller, and differ only in the output
they feed back (Fig.~\ref{fig:aug_closed_loop}).
The controller of each record is known, linear and time invariant, and maps
the position error in actuator coordinates to actuator forces.
For the machine, with controller state $x^{\mathrm{fb}}$, and the model, with
controller state $\hat x^{\mathrm{fb}}$, the loops are
\begin{equation}
\begin{aligned}
 x^{\mathrm{fb}}_{k+1}&=A_{\mathrm{fb}}x^{\mathrm{fb}}_k+B_{\mathrm{fb}}(r_k-y_k),\\
 u_{\mathrm{act},k}&=C_{\mathrm{fb}}x^{\mathrm{fb}}_k+D_{\mathrm{fb}}(r_k-y_k)+u^{\mathrm{ff}}_k,\\
 \hat x^{\mathrm{fb}}_{k+1}&=A_{\mathrm{fb}}\hat x^{\mathrm{fb}}_k+B_{\mathrm{fb}}(r_k-\hat y_k),\\
 \hat u_{\mathrm{act},k}&=C_{\mathrm{fb}}\hat x^{\mathrm{fb}}_k+D_{\mathrm{fb}}(r_k-\hat y_k)+u^{\mathrm{ff}}_k,
\end{aligned}
\label{eq:aug_direct_loop}
\end{equation}
where $(A_{\mathrm{fb}},B_{\mathrm{fb}},C_{\mathrm{fb}},D_{\mathrm{fb}})$ are
the state-space matrices of the controller discretised over $T_s$, and $r_k$,
$u^{\mathrm{ff}}_k\in\R^3$ are the reference and feedforward of the record.

We implement the model loop by subtracting the machine's equations from the
model's, which removes $r$ and $u^{\mathrm{ff}}$, so that only recorded
signals enter.
The result equals the model loop of \eqref{eq:aug_direct_loop} when the
recorded input was produced by this controller from the same output samples
and initial state.
The difference of the controller states needs no estimate: setting it to zero
at the window start places the model's controller in the machine's controller
state, as the encoder places the model state on the recorded trajectory.
Kessels' direct form instead needs the machine's controller state at every
window start, read from the machine or reconstructed from the reference, the
output and the controller~\cite[Remark~5.4]{kessels2025ai}.
\todo{Check for Setup: the training loop uses the controller re-discretised at the model rate, while the records were generated at a higher rate. Run the re-discretised controller on the recorded $r-y$ and compare its feedback force with the recorded one. A rollout with $\hat y=y$ gives $\hat u=u$ by construction and checks nothing.}

\begin{figure}[t]
 \centering
 \includegraphics[width=\columnwidth]{closed_loop_same_reference.pdf}
 \caption{Recorded plant loop (top) and augmented-model loop (bottom),
 driven by one shared reference $r_k$ and feedforward $u_k^{\mathrm{ff}}$,
 with the same feedback controller. The model output is evaluated
 before the input advances its state to the next sample.}
 \label{fig:aug_closed_loop}
\end{figure}

The augmented model has no feedthrough, so each sample is evaluated in a fixed
order: $\hat y_k$ from $\tilde x_k$, then $\hat u_{\mathrm{act},k}$, then the
model and controller states advance.
The controller feedthrough $D_{\mathrm{fb}}$ therefore closes no algebraic
loop, and the loop adds no delay.
The model loop is
\begin{equation}
\begin{aligned}
 x^{c}_{k+1}&=A_{\mathrm{fb}}x^{c}_k+B_{\mathrm{fb}}(y_k-\hat y_k),\qquad x^{c}_\tau=0,\\
 \hat u_{\mathrm{act},k}&=u_{\mathrm{act},k}+C_{\mathrm{fb}}x^{c}_k+D_{\mathrm{fb}}(y_k-\hat y_k),
\end{aligned}
\label{eq:aug_residual_loop}
\end{equation}
where $x^{c}=\hat x^{\mathrm{fb}}-x^{\mathrm{fb}}$,
$\hat u^{\mathrm n}_k=T_u(\hat u_{\mathrm{act},k}-\mu_u)$ replaces
$u^{\mathrm n}_k$ in the model, and $\hat y^{\mathrm n}_{k|\tau}$ denotes the
normalised output of the loop at sample $k\geq\tau$ after the start
\eqref{eq:aug_encoder}.
Model error made before $\tau$ is therefore not carried into the window,
whereas the free-run validation of Section~\ref{sec:aug_closed_loop} simulates
whole records and carries it.

\subsection{Identification problem}\label{sec:aug_closed_loop}
\ifdraft
\begin{itemize}
\item $\theta_{\mathrm{base}}$ is estimated jointly with network and encoder, as in 2026 Eq.~(22); Kessels fixes the first-principles parameters (Eq.~(5.12))
\item Moved from Section~II-C: training coordinates $\xi$, zero at the start; nine positive combinations logarithmic, the signed $m_\Delta$ through \eqref{eq:mdelta_map}, so every iterate satisfies \eqref{eq:lfr_admissibility} without clamp or penalty (D-229); positive-definite inertia, not positive raw masses
\item One display: cost over $\vartheta=(\xi,\theta_{\mathrm{aug}},\theta_{\mathrm{enc}})$ with the model, encoder and loop constraints; every window sample scored; adapted from 2026 Eq.~(22) (closed loop, sample $\tau$ in the encoder); controller fixed
\item A closed-loop fit can hide plant error; Section~VI tests the model under a changed controller
\item Adam, then an L-BFGS polish (adapted from Drenth Sec.~5: unbounded); best validation free-run checkpoint; polish kept only if it improves validation and no meter regresses (D-171, D-172; code); horizons separate, values in Section~V
\item $\mathcal M_{\mathrm U}$ named: zero start, input $\uact$, output $\qact$ (?: notation)
\end{itemize}
\fi

We estimate $\theta_{\mathrm{base}}$ jointly with the network and the encoder,
as Hoekstra et al.\ do~\cite[Eq.~(22)]{hoekstra2026lfr}, whereas Kessels
leaves the first-principles parameters fixed~\cite[Eq.~(5.12)]{kessels2025ai}.
Every iterate must satisfy \eqref{eq:lfr_admissibility}, so that $M(Y)\succ0$
and the LFR stays well posed without a clamp or a penalty.
We therefore train in coordinates $\xi\in\R^{10}$ that are zero at the start
$\theta_{\mathrm{base},0}$.
The nine positive combinations are logarithmic,
$\theta_{\mathrm{base},j}=\theta_{\mathrm{base},0,j}\,e^{\xi_j}$, which keeps
them positive.
The signed $m_\Delta$ is bounded by \eqref{eq:lfr_admissibility} itself:
\begin{equation}
  m_\Delta=\rho\,\frac{2}{L_b}\sqrt{m_\Sigma J_{\mathrm{eff}}}\,\tanh(\eta_0+s\,\xi_\Delta),
  \label{eq:mdelta_map}
\end{equation}
where $\xi_\Delta$ is the component of $\xi$ for $m_\Delta$, $\rho\in(0,1)$ a
margin, and $\eta_0$ and $s$ reproduce $m_{\Delta,0}$ and its sensitivity at
$\xi=0$ (Appendix~\ref{app:lfr-wellposed}).
Every $\xi$ thus gives $M(Y)\succ0$ for all $Y$, which guarantees positive
definite inertia but not positive raw masses.
\todo{Missing: notation decision. The map \eqref{eq:mdelta_map} writes $\theta_{\mathrm{base}}(\xi)$; Section~\ref{sec:obc} and the appendix use the same $\xi$, $\eta_0$, $s$, so the encoder window is now $\chi_\tau$ and the controller state $x^{\mathrm{fb}}$. Candidate: keep this assignment.}

Over the window starts $\mathcal T$ and the window length $n_f$, the trained
parameters $\vartheta=(\xi,\theta_{\mathrm{aug}},\theta_{\mathrm{enc}})$
minimise the normalised closed-loop output error, with the controller fixed.
As in~\cite[Eq.~(22)]{hoekstra2026lfr}, every sample of a window is scored.
We adapt that criterion in two respects: the model is simulated in the loop
of Section~\ref{sec:aug_loop}, and the encoder reads sample $\tau$.
Because the controller acts on the output error, a good closed-loop fit can
mask plant error, so Section~\ref{sec:results} also tests the model under a
changed controller.
The identification problem is
\begin{equation}
\begin{aligned}
 &\min_{\vartheta}\;V(\vartheta)=\frac{1}{|\mathcal T|\,n_f\,n_y}
 \sum_{\tau\in\mathcal T}\sum_{k=\tau}^{\tau+n_f-1}
 \norm{y^{\mathrm n}_{k}-\hat y^{\mathrm n}_{k|\tau}}_2^2\quad\text{s.t.}\\
 &\tilde x^{\mathrm n}_{k+1|\tau}
  =f^{\mathrm n}_{\mathrm{base}}\bigl(\theta_{\mathrm{base}}(\xi),
   \tilde x^{\mathrm n}_{k|\tau},\hat u^{\mathrm n}_{k|\tau}\bigr)
  +f^{\mathrm n}_{\mathrm{aug}}\bigl(\theta_{\mathrm{aug}},
   \hat x^{\mathrm n}_{k|\tau},\hat u^{\mathrm n}_{k|\tau}\bigr),\\
 &\bar x_{k+1|\tau}=g^{\mathrm n}_{\mathrm{aug}}\bigl(\theta_{\mathrm{aug}},
   \hat x^{\mathrm n}_{k|\tau},\hat u^{\mathrm n}_{k|\tau}\bigr),\quad
  \hat y^{\mathrm n}_{k|\tau}=C^{\mathrm n}\tilde x^{\mathrm n}_{k|\tau},\\
 &\hat x^{\mathrm n}_{\tau|\tau}=\psi(\theta_{\mathrm{enc}},\chi_\tau),\quad
  \hat u^{\mathrm n}_{k|\tau}\text{ from \eqref{eq:aug_residual_loop} with }x^{c}_\tau=0,
\end{aligned}
\label{eq:aug_loss}
\end{equation}
where $n_y=3$ and $\theta_{\mathrm{base}}(\xi)$ is the coordinate map above.

We minimise $V$ with Adam on minibatches of windows and then polish with
L-BFGS on a fixed set of windows.
Drenth et al.\ also follow Adam with L-BFGS, in its bounded form
L-BFGS-B~\cite[Sec.~5]{drenth2025lpvlfr}.
Our polish is unbounded.
During Adam, the parameters with the lowest validation error are kept.
This error is the RMS in metres of a closed-loop free run over each whole
validation record, started once from the encoder with $x^{c}=0$, averaged over
the records with their lengths as weights.
The polish is kept only if it lowers this error and no monitored quantity
regresses.
The encoder history $n$, the window length $n_f$ and the free-run horizon are
separate quantities, with values in Section~\ref{sec:setup}.
\todo{Missing: reason for the unbounded polish. Candidate: \eqref{eq:mdelta_map} and the logarithmic coordinates already keep every iterate admissible, so no box bound is needed (own reading; D-171 gives none).}
\todo{Missing: disclosure. The monitored quantities of the polish acceptance include \texttt{combo\_err}, the error of the ten combinations against the true parameters (\texttt{GD/training.py::\_gantry\_meters}), so truth information enters checkpoint acceptance. Candidate: remove \texttt{combo\_err} from the acceptance meters, or state it as an oracle gate in Section~\ref{sec:setup} and report how often it rejected a polish.}

We call $\mathcal M_{\mathrm U}$ the augmented model \eqref{eq:aug_transition}
identified by \eqref{eq:aug_loss} from the zero start \eqref{eq:aug_zero_init},
with $\xi=0$ and the encoder of Section~\ref{sec:aug_encoder}.
Its input is the actuator force $\uact$ and its output the actuator positions
$\qact$.
\todo{Missing: model notation. Candidate: the calligraphic $\mathcal M_{\mathrm U}$ and $\mathcal M_{\mathrm{PS2}}$, since plain $M$ clashes with $M(Y)$ (README open conflict on model names).}

\subsection{Added-state initialisation}\label{sec:aug_ps2}
\ifdraft
\begin{itemize}
\item Problem: from the zero start the added states are zero after one step and $V$ gives $W^a$ and the added-state output rows zero gradient; $V$ can fall through $f_{\mathrm{aug}}$ alone (PS2-METHOD Sec.~2.2; Section~VI tests it)
\item PS2 is this work's, adapted from two-stage regression (Hefny Sec.~2) and its use before backpropagation through time (Downey Sec.~4.2)
\item S1 display: decoder plus $\psi_a$ fitted to the current model's normalised residual over $M$ samples; trains $\theta_{\mathrm p}$, $\theta^a_{\mathrm{enc}}$ only, own optimiser, after every Adam update; calibration records withheld; fallback if unpredictable (code)
\item S2 display: once, one-step regression of $g^{\mathrm n}_{\mathrm{aug}}$ on frozen encoder pairs; physical output rows unchanged, hidden layers move (code)
\item S3: decoder discarded, \eqref{eq:aug_loss} continues unchanged; deployed model, criterion and selection unchanged; added states latent; $\mathcal M_{\mathrm{PS2}}$ named
\end{itemize}
\fi

The zero start \eqref{eq:aug_zero_init} leaves the added states without a
role.
They are zero after the first step, and the output depends on them only
through the zero output layer, so the gradient of $V$ with respect to $W^a$
and to the added-state rows of $W_L$ is zero at the start.
Minimising $V$ can then lower the error through $f^{\mathrm n}_{\mathrm{aug}}$
alone, with static corrections or real lags, and leave the omitted dynamics
unrepresented.
Section~\ref{sec:results} tests whether it does.
We therefore propose PS2, which gives the added states a predictive role from
data before the network reads them and then hands over to \eqref{eq:aug_loss}.
It adapts the two-stage regression of predictive-state
learning~\cite[Sec.~2]{hefny2015supervised} and its use as an initialisation
before backpropagation through time~\cite[Sec.~4.2]{downey2017psrnn}.
\todo{Missing: the evidence for plain training failing in this way is from the previous campaign (\texttt{PS2-METHOD.md} Sec.~2.2, detuned runs, older data). Candidate: keep the sentence conditional, as written, until \texttt{THESIS-RESULTS.md} step~2 reports it.}

Stage S1 trains the encoder's added-state outputs to predict the residual of
the model being trained.
During the first updates of \eqref{eq:aug_loss}, one S1 update with its own
optimiser follows every Adam update.
S1 uses windows from the training records except a calibration subset.
Its target is the current model's closed-loop residual over the next $M$
samples, with $M=n+1$ the length of the encoder window, so that the past and
future horizons are equal, a convention of subspace identification.
Unlike the first-stage regression of~\cite{hefny2015supervised}, which
predicts features of future observations, S1 predicts this residual through a
temporary decoder that reads only $\bar x_\tau$:
\begin{equation}
\begin{aligned}
 e_\tau&=\frac{1}{\sigma_e}\col\bigl(y^{\mathrm n}_{k}-\hat y^{\mathrm n}_{k|\tau}\bigr)_{k=\tau}^{\tau+M-1},\\
 &\min_{\theta_{\mathrm p},\,\theta^a_{\mathrm{enc}}}\;
 \frac{1}{|\mathcal B|\,M\,n_y}\sum_{\tau\in\mathcal B}
 \norm{d\bigl(\theta_{\mathrm p},\psi_a(\theta_{\mathrm{enc}},\chi_\tau)\bigr)-e_\tau}_2^2,
\end{aligned}
\label{eq:ps2_s1}
\end{equation}
where $\hat y^{\mathrm n}_{k|\tau}$ is the output of the loop of
Section~\ref{sec:aug_loop} at the current $\vartheta$, held fixed in the
gradient, $\sigma_e$ the RMS of the residuals over the minibatch
$\mathcal B$, and
$d:\R^{n_{\mathrm p}}\times\R^{n_{\bar x}}\to\R^{Mn_y}$ a decoder with
parameters $\theta_{\mathrm p}\in\R^{n_{\mathrm p}}$, one tanh layer of the
width of $\phi_{\mathrm{aug}}$ and a zero output layer.

S1 ends when the fraction of the residual that the decoder leaves unexplained
on the calibration records stops falling, or after a maximum number of
updates.
If that fraction is then too large, S2 is skipped and training continues
with the S1 changes kept.
Otherwise, stage S2 fits the added-state transition so that it propagates
these states.
Once, at the end of S1, we encode window pairs that start at samples $k_i$
and $k_i+1$ anywhere in the fit records, freeze them, and fit
$g^{\mathrm n}_{\mathrm{aug}}$ by one-step regression with a fresh optimiser,
stopped on pairs from the calibration records.
Unlike the linear operator of the second stage
in~\cite{hefny2015supervised}, the regressed map is the nonlinear transition
of the deployed model:
\begin{equation}
 \min_{\theta_{\mathrm{aug}}}\;
 \frac{\sum_{i}\norm{g^{\mathrm n}_{\mathrm{aug}}\bigl(\theta_{\mathrm{aug}},
   \psi(\theta_{\mathrm{enc}},\chi_{k_i}),u^{\mathrm n}_{k_i}\bigr)
   -\psi_a(\theta_{\mathrm{enc}},\chi_{k_i+1})}_2^2}
 {\sum_{j=1}^{n_{\bar x}}\sum_{i}\bigl(\psi_{a,j}(\theta_{\mathrm{enc}},\chi_{k_i+1})-\bar\psi_{a,j}\bigr)^2},
 \label{eq:ps2_s2}
\end{equation}
where $i$ runs over a minibatch of frozen pairs, $u^{\mathrm n}_{k_i}$ is the
recorded input, $\theta_{\mathrm{enc}}$ is held at its value at the end of S1,
and $\bar\psi_{a,j}$ is the minibatch mean of the $j$th target.
The loss reaches $\phi_{\mathrm{aug}}$ only through its added-state rows, so
S2 leaves the physical rows of $W_L$ unchanged, while the shared hidden layers
move.
\todo{Missing: reason for regressing on the recorded input $u^{\mathrm n}_{k_i}$ rather than the closed-loop input of \eqref{eq:aug_residual_loop}. Candidate: at the window start $x^{c}=0$, so the two differ only by $D_{\mathrm{fb}}(y-\hat y)$ (own reading; \texttt{ps2\_opening.py::\_fit\_map} uses the recorded input; no decision entry states a reason).}

In stage S3 the decoder is discarded and the minimisation of
\eqref{eq:aug_loss} continues from the modified parameters, followed by the
polish and selection of Section~\ref{sec:aug_closed_loop}.
PS2 adds nothing to the deployed model, the criterion or the selection rule.
The added states remain latent, since PS2 assigns them no physical
coordinates.
An invertible transformation of $\bar x$, absorbed by $\phi_{\mathrm{aug}}$,
leaves the input-output behaviour unchanged.
We call $\mathcal M_{\mathrm{PS2}}$ the model \eqref{eq:aug_transition}
identified by \eqref{eq:aug_loss} with this opening, from the same start as
$\mathcal M_{\mathrm U}$.
\todo{Missing: basis of the PS2 switch and stop rules (S1 plateau and cap, screen threshold, S2 plateau and cap, snapshot sizes, calibration split as the last record of each class with three or more records). Candidate: rows of the hyperparameter table in Section~\ref{sec:setup}, all HEURISTIC (\texttt{PS2Spec}; \texttt{PS2-METHOD.md} Sec.~5), noting that the S1 cap, not the plateau, ended S1 in every run so far.}
